% GENERATED FILE. Edit canonical lessons, then run npm run generate:books.
% canonical-source-hash: fnv1a64:66db42ab7e3e4a10
% canonical-lessons: TA-C320-kotiyerru, TA-C320-anivakuppu, TA-C320-toranam, TA-C320-patakai, TA-C320-paraicarru

\chapter{Flags and Parades}
\label{ch:ta-flags-and-parades}
\hlchaptermodality{320}

\begin{chapteropening}
\textbf{By the end of this chapter:} \emph{I can describe a flag-raising and a parade.}

Complete the last lesson of chapter 320: I can describe a flag-raising and a parade.
\end{chapteropening}

\section[\texorpdfstring{\ta{கொடியேற்று} (koṭiyēṟṟu)}{கொடியேற்று (koṭiyēṟṟu)}]{\ta{கொடியேற்று} (koṭiyēṟṟu) — to hoist a flag}

\label{lesson:TA-C320-kotiyerru}



\begin{quote}
Before the new one: say the Tamil for a festive look, then the Tamil for an invitation card.
\end{quote}

\subsection*{You'll want to know: \ta{கொடியேற்று}}
\textbf{\ta{கொடியேற்று}} — \emph{koṭiyēṟṟu} — \textquotedblleft{}to hoist a flag\textquotedblright{}.

\begin{quote}
\ta{தலைமை} \ta{ஆசிரியர்} \ta{காலையில்} \ta{கொடியேற்றினார்}.
\end{quote}

\begin{quote}
\emph{talaimai āciriyar kālaiyil koṭiyēṟṟiṉār.}
\end{quote}

\begin{quote}
The headmaster hoisted the flag in the morning.
\end{quote}

Past (I …): \ta{கொடியேற்றினேன்}, \emph{koṭiyēṟṟiṉēṉ}.

\ta{கொடி}, flag + \ta{ஏற்று}, to raise.

\begin{cousinweb}
\textbf{\ta{கொடியேற்று}} (\emph{koṭiyēṟṟu}) is a Tamil compound built from \ta{கொடி}, \textquotedblleft{}banner, flag\textquotedblright{} (DEDR 2049), and \ta{ஏற்று}, \textquotedblleft{}to lift up, raise\textquotedblright{} (DEDR 916).
\end{cousinweb}

\begin{grammarlens}[title={What you've built}]
One more word for telling what happened.
\end{grammarlens}

\subsection*{Your turn}
\begin{itemize}
\raggedright
  \item \emph{Say it:} \emph{koṭiyēṟṟu}
  \item \emph{Say it:} \emph{koṭiyēṟṟu}, once more
  \item \emph{Say it:} \emph{talaimai āciriyar kālaiyil koṭiyēṟṟiṉār.} — The headmaster hoisted the flag in the morning.
  \item \emph{Recall:} say the Tamil for a festive look, then the Tamil for an invitation card, then say \emph{koṭiyēṟṟu} again
\end{itemize}

\begin{tcolorbox}[breakable,colback=teal!4,colframe=teal!35!black,title={Before you move on}]
What does \ta{கொடியேற்று} mean? (\textquotedblleft{}To hoist a flag\textquotedblright{}.) Say it once more.
\end{tcolorbox}
\section[\texorpdfstring{\ta{அணிவகுப்பு} (aṇivakuppu)}{அணிவகுப்பு (aṇivakuppu)}]{\ta{அணிவகுப்பு} (aṇivakuppu) — a parade; a march past}

\label{lesson:TA-C320-anivakuppu}



\begin{quote}
Before the new one: say the Tamil for an invitation card, then the Tamil for to hoist a flag.
\end{quote}

\subsection*{You'll want to know: \ta{அணிவகுப்பு}}
\textbf{\ta{அணிவகுப்பு}} — \emph{aṇivakuppu} — \textquotedblleft{}a parade; a march past\textquotedblright{}.

\begin{quote}
\ta{மாணவர்கள்} \ta{அணிவகுப்பில்} \ta{கலந்துகொண்டார்கள்}.
\end{quote}

\begin{quote}
\emph{māṇavarkaḷ aṇivakuppil kalantukoṇṭārkaḷ.}
\end{quote}

\begin{quote}
The students took part in the parade.
\end{quote}

\ta{அணி}, a row + \ta{வகு}, to divide: ranks marching in order.

\begin{cousinweb}
\textbf{\ta{அணிவகுப்பு}} (\emph{aṇivakuppu}) is a Tamil compound built from \ta{அணி}, \textquotedblleft{}array of an army, row\textquotedblright{} (DEDR 117, which asks whether it comes from Sanskrit \emph{anīka}), and \ta{வகுப்பு}, \textquotedblleft{}division\textquotedblright{}, from \ta{வகு}, \textquotedblleft{}to divide\textquotedblright{} (DEDR 5202).
\end{cousinweb}

\begin{grammarlens}[title={What you've built}]
One more word for telling what happened.
\end{grammarlens}

\subsection*{Your turn}
\begin{itemize}
\raggedright
  \item \emph{Say it:} \emph{aṇivakuppu}
  \item \emph{Say it:} \emph{aṇivakuppu}, once more
  \item \emph{Say it:} \emph{māṇavarkaḷ aṇivakuppil kalantukoṇṭārkaḷ.} — The students took part in the parade.
  \item \emph{Recall:} say the Tamil for an invitation card, then the Tamil for to hoist a flag, then say \emph{aṇivakuppu} again
\end{itemize}

\begin{tcolorbox}[breakable,colback=teal!4,colframe=teal!35!black,title={Before you move on}]
What does \ta{அணிவகுப்பு} mean? (\textquotedblleft{}A parade; a march past\textquotedblright{}.) Say it once more.
\end{tcolorbox}
\section[\texorpdfstring{\ta{தோரணம்} (tōraṇam)}{தோரணம் (tōraṇam)}]{\ta{தோரணம்} (tōraṇam) — a festoon}

\label{lesson:TA-C320-toranam}



\begin{quote}
Before the new one: say the Tamil for to hoist a flag, then the Tamil for a parade.
\end{quote}

\subsection*{You'll want to know: \ta{தோரணம்}}
\textbf{\ta{தோரணம்}} — \emph{tōraṇam} — \textquotedblleft{}a festoon\textquotedblright{}.

\begin{quote}
\ta{தெருவில்} \ta{வண்ணத்} \ta{தோரணங்கள்} \ta{கட்டியிருந்தார்கள்}.
\end{quote}

\begin{quote}
\emph{teruvil vaṇṇat tōraṇaṅkaḷ kaṭṭiyiruntārkaḷ.}
\end{quote}

\begin{quote}
They had hung colourful festoons across the street.
\end{quote}

Plural: \ta{தோரணங்கள்}, \emph{tōraṇaṅkaḷ}.

Plural \ta{தோரணங்கள்}.

\begin{cousinweb}
\textbf{\ta{தோரணம்}} (\emph{tōraṇam}) is borrowed from Sanskrit \emph{tōraṇa}, \textquotedblleft{}an arched or festooned doorway\textquotedblright{} (CDIAL 5976).
\end{cousinweb}

\begin{grammarlens}[title={What you've built}]
One more word for telling what happened.
\end{grammarlens}

\subsection*{Your turn}
\begin{itemize}
\raggedright
  \item \emph{Say it:} \emph{tōraṇam}
  \item \emph{Say it:} \emph{tōraṇam}, once more
  \item \emph{Say it:} \emph{teruvil vaṇṇat tōraṇaṅkaḷ kaṭṭiyiruntārkaḷ.} — They had hung colourful festoons across the street.
  \item \emph{Recall:} say the Tamil for to hoist a flag, then the Tamil for a parade, then say \emph{tōraṇam} again
\end{itemize}

\begin{tcolorbox}[breakable,colback=teal!4,colframe=teal!35!black,title={Before you move on}]
What does \ta{தோரணம்} mean? (\textquotedblleft{}A festoon\textquotedblright{}.) Say it once more.
\end{tcolorbox}
\section[\texorpdfstring{\ta{பதாகை} (patākai)}{பதாகை (patākai)}]{\ta{பதாகை} (patākai) — a banner}

\label{lesson:TA-C320-patakai}



\begin{quote}
Before the new one: say the Tamil for a parade, then the Tamil for a festoon.
\end{quote}

\subsection*{You'll want to know: \ta{பதாகை}}
\textbf{\ta{பதாகை}} — \emph{patākai} — \textquotedblleft{}a banner\textquotedblright{}.

\begin{quote}
\ta{மாணவர்கள்} \ta{பெரிய} \ta{பதாகைகளைப்} \ta{பிடித்துக்கொண்டு} \ta{நடந்தார்கள்}.
\end{quote}

\begin{quote}
\emph{māṇavarkaḷ periya patākaikaḷaip piṭittukkoṇṭu naṭantārkaḷ.}
\end{quote}

\begin{quote}
The students marched holding big banners.
\end{quote}

Plural: \ta{பதாகைகள்}, \emph{patākaikaḷ}.

Carried in parades and protests; the plural is \ta{பதாகைகள்}.

\begin{cousinweb}
\textbf{\ta{பதாகை}} (\emph{patākai}) is borrowed from Sanskrit \emph{patākā}, \textquotedblleft{}a flag\textquotedblright{} (CDIAL 7726); the Sanskrit final -ā becomes Tamil -\ta{ஐ}.
\end{cousinweb}

\begin{grammarlens}[title={What you've built}]
One more word for telling what happened.
\end{grammarlens}

\subsection*{Your turn}
\begin{itemize}
\raggedright
  \item \emph{Say it:} \emph{patākai}
  \item \emph{Say it:} \emph{patākai}, once more
  \item \emph{Say it:} \emph{māṇavarkaḷ periya patākaikaḷaip piṭittukkoṇṭu naṭantārkaḷ.} — The students marched holding big banners.
  \item \emph{Recall:} say the Tamil for a parade, then the Tamil for a festoon, then say \emph{patākai} again
\end{itemize}

\begin{tcolorbox}[breakable,colback=teal!4,colframe=teal!35!black,title={Before you move on}]
What does \ta{பதாகை} mean? (\textquotedblleft{}A banner\textquotedblright{}.) Say it once more.
\end{tcolorbox}
\section[\texorpdfstring{\ta{பறைசாற்று} (paṟaicāṟṟu)}{பறைசாற்று (paṟaicāṟṟu)}]{\ta{பறைசாற்று} (paṟaicāṟṟu) — to proclaim}

\label{lesson:TA-C320-paraicarru}



\begin{quote}
Before the new one: say the Tamil for a festoon, then the Tamil for a banner.
\end{quote}

\subsection*{You'll want to know: \ta{பறைசாற்று}}
\textbf{\ta{பறைசாற்று}} — \emph{paṟaicāṟṟu} — \textquotedblleft{}to proclaim\textquotedblright{}.

\begin{quote}
\ta{இந்தக்} \ta{கோயில்} \ta{பழைய} \ta{கலையின்} \ta{பெருமையைப்} \ta{பறைசாற்றுகிறது}.
\end{quote}

\begin{quote}
\emph{intak kōyil paḻaiya kalaiyiṉ perumaiyaip paṟaicāṟṟukiṟatu.}
\end{quote}

\begin{quote}
This temple proclaims the glory of the old art.
\end{quote}

Past (I …): \ta{பறைசாற்றினேன்}, \emph{paṟaicāṟṟiṉēṉ}.

\ta{பறை}, drum + \ta{சாற்று}, to announce: once by drum, now any loud proof.

\begin{cousinweb}
\textbf{\ta{பறைசாற்று}} (\emph{paṟaicāṟṟu}) is a Tamil compound built from \ta{பறை}, \textquotedblleft{}drum\textquotedblright{} (DEDR 4032), and \ta{சாற்று}, \textquotedblleft{}to publish, announce, beat as a drum\textquotedblright{} (DEDR 2486).
\end{cousinweb}

\begin{grammarlens}[title={What you've built}]
One more word for telling what happened.
\end{grammarlens}

\subsection*{Your turn}
\begin{itemize}
\raggedright
  \item \emph{Say it:} \emph{paṟaicāṟṟu}
  \item \emph{Say it:} \emph{paṟaicāṟṟu}, once more
  \item \emph{Say it:} \emph{intak kōyil paḻaiya kalaiyiṉ perumaiyaip paṟaicāṟṟukiṟatu.} — This temple proclaims the glory of the old art.
  \item \emph{Recall:} say the Tamil for a festoon, then the Tamil for a banner, then say \emph{paṟaicāṟṟu} again
\end{itemize}

\begin{tcolorbox}[breakable,colback=teal!4,colframe=teal!35!black,title={Before you move on}]
What does \ta{பறைசாற்று} mean? (\textquotedblleft{}To proclaim\textquotedblright{}.) Say it once more.
\end{tcolorbox}
